\chapter{Modelos Bayesianos e de Fronteira}

\section{Regressão linear bayesiana}

\texttt{bayes\_lm} implementa regressão linear bayesiana com priori conjugada
Normal-Inversa-Gamma. Com uma priori fraca, recupera as médias
posteriori do OLS, mas também reporta desvios padrão posteriores,
intervalos credíveis e a probabilidade de cada coeficiente ser positivo.

\begin{lstlisting}[caption={Regressão linear bayesiana}]
set_seed(42)
let n = 200
generate df x1 = rnormal(n)
generate df x2 = rnormal(n)
generate df y = 1 + 2*x1 - 1.5*x2 + rnormal(n, 0, 2)

let m_bayes = bayes_lm(y ~ x1 + x2, df)
print(m_bayes)
\end{lstlisting}

A média posterior de cada coeficiente é a média ponderada entre a média
a priori e a estimativa OLS, com mais peso nos dados à medida que a
amostra cresce. O intervalo credível de 95\% é reportado para cada
regressor.

\section{Análise de fronteira estocástica}

Modelos de fronteira estocástica separam ruído estatístico de
ineficiência técnica ou de custo. A fronteira de produção é

\[
y = X\beta + v - u, \qquad u \ge 0
\]

onde $v$ é o ruído e $u$ é a ineficiência. \texttt{sfa\_production} estima
$\beta$, $\sigma_v$, $\sigma_u$ e a eficiência técnica média.

\begin{lstlisting}[caption={Fronteira de produção estocástica}]
let n = 300
generate df x1 = rnormal(n)
generate df x2 = rnormal(n)
generate df v  = rnormal(n, 0, 0.3)
generate df u  = abs(rnormal(n, 0, 0.5))
generate df y  = 1 + 0.8*x1 + 0.6*x2 + v - u

let m_sfa = sfa_production(y ~ x1 + x2, df)
print(m_sfa)
\end{lstlisting}

Para uma fronteira de custo, o termo de ineficiência é somado em vez de
subtraído:

\begin{lstlisting}[caption={Fronteira de custo estocástica}]
let m_sfac = sfa_cost(y ~ x1 + x2, df)
print(m_sfac)
\end{lstlisting}

\section{Fronteira estocástica bayesiana}

\texttt{bayes\_sfa\_production} e \texttt{bayes\_sfa\_cost} implementam um amostrador de
Gibbs com prioris conjugadas fracas. Eles reportam médias posteriores,
desvios padrão e intervalos credíveis de 95\% para os coeficientes de
regressão e os parâmetros de variância.

\begin{lstlisting}[caption={SFA bayesiano}]
let m_bsfa = bayes_sfa_production(y ~ x1 + x2, df, burn=500, draws=1000)
print(m_bsfa)
\end{lstlisting}

São mais robustos em amostras pequenas que SFA por máxima
verossimilhança, mas a convergência deve ser verificada rodando a
cadeira com diferentes valores iniciais e comparando as posteriores.

\section{Quando usar cada um}

\begin{itemize}
\item \textbf{Inferência padrão com incerteza:} \texttt{bayes\_lm}
\item \textbf{Eficiência de produção:} \texttt{sfa\_production}
\item \textbf{Eficiência de custo:} \texttt{sfa\_cost}
\item \textbf{SFA em amostras pequenas ou posteriori completa:} \texttt{bayes\_sfa\_production}
\end{itemize}
